\documentclass[11pt]{article}

\usepackage[T1]{fontenc}
\usepackage[utf8]{inputenc}
\usepackage{lmodern}
\usepackage[a4paper,margin=29mm]{geometry}
\usepackage{microtype}
\usepackage{amsmath,amssymb,amsthm,mathtools}
\usepackage{booktabs}
\usepackage{array}
\usepackage{enumitem}
\usepackage{xcolor}
\usepackage{hyperref}
\usepackage[nameinlink,noabbrev]{cleveref}

\hypersetup{
  colorlinks=true,
  linkcolor=blue!45!black,
  citecolor=green!35!black,
  urlcolor=blue!55!black,
  pdftitle={Rank-two saturation, algebraic-orbit collapse, and Cartier-state growth},
  pdfauthor={OpenAI, prepared for Vladimir Reshetnikov}
}

\setlist[itemize]{leftmargin=1.7em,itemsep=0.25em,topsep=0.35em}
\setlist[enumerate]{leftmargin=1.9em,itemsep=0.25em,topsep=0.35em}
\allowdisplaybreaks
\emergencystretch=3em

\newtheorem{theorem}{Theorem}[section]
\newtheorem{proposition}[theorem]{Proposition}
\newtheorem{lemma}[theorem]{Lemma}
\newtheorem{corollary}[theorem]{Corollary}
\newtheorem{conjecture}[theorem]{Conjectural bridge}
\theoremstyle{definition}
\newtheorem{definition}[theorem]{Definition}
\newtheorem{remark}[theorem]{Remark}
\newtheorem{problem}[theorem]{Target problem}

\newcommand{\Q}{\mathbb{Q}}
\newcommand{\Z}{\mathbb{Z}}
\newcommand{\N}{\mathbb{N}}
\newcommand{\R}{\mathbb{R}}
\newcommand{\F}{\mathbb{F}}
\newcommand{\Qbar}{\overline{\mathbb{Q}}}
\newcommand{\rank}{\operatorname{rank}}
\newcommand{\ord}{\operatorname{ord}}
\newcommand{\rad}{\operatorname{rad}}
\newcommand{\Span}{\operatorname{span}}
\newcommand{\supp}{\operatorname{supp}}
\newcommand{\val}{\boldsymbol{v}}
\newcommand{\smooth}{\mathcal{S}}
\newcommand{\Gammao}{\Gamma_0}
\newcommand{\Gammai}{\Gamma_1}
\newcommand{\Orb}{\mathcal{O}}
\newcommand{\Alg}{\mathcal{A}}
\newcommand{\Cart}{\Lambda}

\title{\bfseries Rank-Two Saturation, Algebraic-Orbit Collapse,\\
and Cartier-State Growth\\[0.45em]
\large New progress on the Alaoglu--Erd\H{o}s exponent problem}
\author{Research report prepared for Vladimir Reshetnikov}
\date{23 August 2026}

\begin{document}
\maketitle

\begin{abstract}
Let $x\in\R$ and suppose
\[
        M=2^x\in\Z,\qquad N=3^x\in\Z.
\]
The Alaoglu--Erd\H{o}s conjecture asserts that $x\in\Z$.  This report does not
claim a complete proof.  It develops a set of deductions that are independent
of the determinant, generic interpolation, density, and broad $p$-adic routes
already developed in the supplied unified report.

The main unconditional result is an \emph{algebraic-orbit filtration}.  Under
a hypothetical nonintegral solution, the positive algebraic numbers whose
$x$-th powers remain algebraic form exactly the rank-two radical generated by
$2$ and $3$.  The subgroup whose first two iterates
$\alpha,\alpha^x,\alpha^{x^2}$ are algebraic has rank at most one, while no
nontrivial positive algebraic number can have four consecutive algebraic
iterates
$\alpha,\alpha^x,\alpha^{x^2},\alpha^{x^3}$.  This uses the Six Exponentials
Theorem twice, followed by a finite-rank orbit argument.  It yields a sharp
rank-$3$/rank-$4$ dichotomy for the multiplicative group generated by
$2,3,M,N$, and it forces transcendence by the third iterate along the unique
return direction in the rank-$3$ case.

A second result gives complete \emph{additive isolation}: except for the four
consecutive $\{2,3\}$-smooth pairs
$(1,2),(2,3),(3,4),(8,9)$, every successor of a positive $\{2,3\}$-smooth
integer has transcendental $x$-th power.  Thus a counterexample would be
multiplicatively closed but maximally hostile to even the smallest additive
propagation.

Finally, a weighted smooth-support series is introduced,
\[
  H_{M,N}(z)=\sum_{a,b\geq0}M^aN^b z^{2^a3^b}.
\]
It obeys an exact mixed $2$--$3$ Mahler equation.  We prove exponential lower
bounds for the number of binary and ternary Cartier states of the unweighted
support series, and classify completely when $H_{M,N}$ becomes algebraic
after reduction modulo any prime.  These results expose both the power and the
present limitation of the automatic/Mahler route: finite-characteristic
behavior detects the prime-divisibility geometry of $M,N$, whereas the
conjecture is controlled by an archimedean equality of two pole abscissae.  A
successful proof needs a genuine bridge between those two layers.
\end{abstract}

\tableofcontents

\section{Scope, status, and notation}

The purpose of this report is to add a nonoverlapping layer to the supplied
unified report.  In particular, it does not repeat the standard reduction to a
$2\times2$ logarithmic determinant, the usual Four Exponentials implication,
generic interpolation determinants, broad smooth-number density estimates,
or exploratory $p$-adic substitutions.  Those routes remain relevant, but
repeating them would obscure the genuinely new deductions below.

All statements labelled theorem, proposition, lemma, or corollary are proved
here from standard named results.  Statements labelled target problem or
conjectural bridge are research objectives, not established facts.  The
reported work of other systems is not used as a premise anywhere in the
mathematics.

Throughout, a \emph{counterexample} means a real $x\notin\Z$ for which
\[
        M:=2^x\in\Z_{>0},\qquad N:=3^x\in\Z_{>0}.
\]
The case $x=0$ is the trivial solution $M=N=1$.  For a finitely generated
multiplicative group $G$ of positive algebraic numbers, write
\[
 \rad_{\Qbar}(G)
 :=\{\alpha\in\Qbar_{>0}:\alpha^k\in G
       \text{ for some }k\geq1\}.
\]
Its multiplicative rank is the dimension of $G\otimes_{\Z}\Q$.  We use
\[
  \Gammao:=\langle2,3\rangle\subset\Q_{>0},\qquad
  \Gammai:=\langle M,N\rangle\subset\Q_{>0},
\]
and
\[
  \smooth:=\Gammao\cap\N=\{2^a3^b:a,b\in\Z_{\geq0}\}.
\]
Positive real powers are always taken with the ordinary real logarithm, so no
branch choices occur.

\section{Elementary normalization and the ordered-semigroup form}

We begin with a compact normalization package.  The individual assertions are
simple, but keeping them explicit prevents later arguments from silently using
more than is known.

\begin{proposition}[Algebraic exponents are already excluded]
\label{prop:alg-exponent}
If $2^x$ and $3^x$ are integers and $x$ is algebraic, then $x\in\Z$.
Consequently every counterexample is positive and transcendental.
\end{proposition}

\begin{proof}
If $x<0$, then $0<2^x<1$, so $2^x$ is not an integer.  Suppose first that
$x=a/b\in\Q$ in lowest terms, with $b>0$.  From $2^a=M^b$, unique
factorization gives $M=2^c$ and $a=bc$.  Hence $b\mid a$, so $b=1$.

If $x$ is algebraic and irrational, the Gelfond--Schneider theorem~\cite{Gelfond,Schneider} says that
$2^x$ is transcendental, contradicting $M\in\Z$.  Thus an algebraic solution
is rational and therefore integral.  A nonintegral solution must consequently
be positive and transcendental.
\end{proof}

\begin{proposition}[Free ordered-semigroup realization]
\label{prop:ordered-semigroup}
For a counterexample, $M$ and $N$ are multiplicatively independent, and
\[
 \Phi_x:\Gammao\longrightarrow\Gammai,
 \qquad \Phi_x(2^a3^b)=M^aN^b,
\]
is an injective multiplicative isomorphism.  It is strictly order preserving:
for $u,v\in\Gammao$,
\[
       u<v \quad\Longleftrightarrow\quad \Phi_x(u)<\Phi_x(v).
\]
Moreover $\log \Phi_x(u)=x\log u$.
\end{proposition}

\begin{proof}
If $M^aN^b=1$ for integers $a,b$, then
$2^{ax}3^{bx}=1$.  Since $x>0$, this gives $2^a3^b=1$, hence $a=b=0$.
Thus $M,N$ are multiplicatively independent and the displayed map is an
isomorphism of free abelian groups.  The logarithmic identity follows directly
from $M=2^x$ and $N=3^x$; strict order preservation follows because $x>0$.
\end{proof}

This gives a useful exact restatement: the problem asks whether the ordered
multiplicative group generated by $2,3$ has any integral realization other
than the integral power realizations
\[
          (2,3)\longmapsto(2^k,3^k),\qquad k\in\Z_{\geq0}.
\]
The word ``ordered'' is essential.  As abstract free abelian groups there are,
of course, infinitely many embeddings into $\Q_{>0}$.

\begin{definition}[Perfect-power content]
For an integer $A>1$, put
\[
  \operatorname{cont}(A):=\gcd\{v_p(A):p\mid A\}.
\]
Thus $A$ is a perfect $d$-th power exactly when $d\mid\operatorname{cont}(A)$.
\end{definition}

\begin{proposition}[Primitive core]
\label{prop:primitive-core}
Let $x$ be a counterexample and let
\[
 d=\gcd\bigl(\operatorname{cont}(M),\operatorname{cont}(N)\bigr).
\]
Then $M=M_0^d$ and $N=N_0^d$ for integers $M_0,N_0$, and $x_0=x/d$ is again a
counterexample:
\[
         2^{x_0}=M_0,
         \qquad 3^{x_0}=N_0.
\]
After dividing by the maximal such $d$, one may therefore assume that the
combined prime-exponent content of $M$ and $N$ is $1$.
\end{proposition}

\begin{proof}
The definition of $d$ gives integer positive $d$-th roots $M_0,N_0$.  Taking
the positive real $d$-th root of $2^x=M_0^d$ gives $2^{x/d}=M_0$, and similarly
for $3$.  If $x/d$ were integral, then $x$ would be integral as well, contrary
to hypothesis.  Maximality gives the final assertion.
\end{proof}

This is a genuine descent, although it only removes a common perfect-power
factor from the two outputs.  Separate perfect-power exponents of $M$ and $N$
do not by themselves give a common descent.

\section{The algebraic-power domain is exactly rank two}

The next result is the first major structural restriction.  It uses the
classical Six Exponentials Theorem in its symmetric form: if
$u_1,u_2$ are linearly independent over $\Q$ and $v_1,v_2,v_3$ are linearly
independent over $\Q$, then at least one of the six numbers
$\exp(u_i v_j)$ is transcendental.

\begin{definition}[Algebraic-power domain]
For $t\in\R$, define
\[
 \Alg(t):=\{\alpha\in\Qbar_{>0}:\alpha^t\in\Qbar\}.
\]
This is a radical multiplicative subgroup of $\Qbar_{>0}$.
\end{definition}

\begin{lemma}[Six-exponential rank bound]
\label{lem:six-rank}
If $t\notin\Q$, then $\Alg(t)$ has multiplicative rank at most $2$.
\end{lemma}

\begin{proof}
If $\alpha_1,\alpha_2,\alpha_3\in\Alg(t)$ were multiplicatively independent,
then $\log\alpha_1,\log\alpha_2,\log\alpha_3$ would be linearly independent
over $\Q$, as would $1,t$.  All six numbers
\[
  \exp(1\cdot\log\alpha_j)=\alpha_j,
  \qquad
  \exp(t\log\alpha_j)=\alpha_j^t
\]
would be algebraic, contradicting the Six Exponentials Theorem.
\end{proof}

\begin{theorem}[Exact saturation of the algebraic-power domain]
\label{thm:saturation}
Assume that $x$ is a counterexample.  Then
\[
       \Alg(x)=\rad_{\Qbar}(\Gammao).
\]
Furthermore the positive-power map $\alpha\mapsto\alpha^x$ is an isomorphism
\[
   \rad_{\Qbar}(\Gammao)
      \xrightarrow{\ \sim\ }
   \rad_{\Qbar}(\Gammai).
\]
In particular,
\[
 \{q\in\Q_{>0}:q^x\in\Qbar\}=\Gammao.
\]
Thus, for every positive rational $q$ having a prime factor other than $2$ or
$3$, the number $q^x$ is transcendental.
\end{theorem}

\begin{proof}
By hypothesis $2,3\in\Alg(x)$, and they are multiplicatively independent.
By \cref{lem:six-rank}, no element $\alpha\in\Alg(x)$ can be multiplicatively
independent from both $2$ and $3$.  Hence there exist integers $k\geq1$ and
$a,b$ such that
\[
          \alpha^k=2^a3^b,
\]
so $\alpha\in\rad_{\Qbar}(\Gammao)$.

Conversely, if $\alpha^k=2^a3^b$, then
\[
       (\alpha^x)^k=(\alpha^k)^x=M^aN^b\in\Qbar,
\]
so $\alpha^x$ is algebraic.  This proves the first equality.

The same computation shows that the image lies in
$\rad_{\Qbar}(\Gammai)$.  Conversely, if $\beta^k=M^aN^b$, let $\alpha$ be the
positive $k$-th root of $2^a3^b$.  Then $\alpha\in\rad_{\Qbar}(\Gammao)$ and,
by uniqueness of the positive $k$-th root, $\alpha^x=\beta$.  This proves
surjectivity and hence the displayed isomorphism.

Finally, if a rational $q$ belongs to $\rad_{\Qbar}(\Gammao)$, then
$q^k=2^a3^b$ for some integers.  Taking ordinary prime valuations shows that
$q=2^r3^s$ with $r,s\in\Z$.  The rational part of the radical is therefore
exactly $\Gammao$.
\end{proof}

\begin{corollary}[One-more-base principle]
\label{cor:one-more-base}
Under a hypothetical counterexample, proving the algebraicity of $q^x$ for
\emph{one} rational $q\notin\Gammao$ would finish the conjecture.  In
particular, any one of
\[
          5^x,\quad 7^x,\quad 10^x,\quad 11^x
\]
being algebraic is already impossible.
\end{corollary}

This sharpens the familiar slogan that three multiplicatively independent
bases suffice.  It says more: in a counterexample, the two known bases
$2,3$ exhaust the entire algebraic-power domain, including all algebraic
radicals.  Every successful strategy must therefore manufacture a genuinely
new multiplicative direction; taking more products and quotients of the known
integer values can never do so.

\section{Algebraic-orbit filtration and collapse by the third iterate}

The saturation theorem admits an unexpectedly rigid iteration.  This is the
strongest new deduction in the report.

\begin{definition}[Algebraic-orbit filtration]
For an integer $r\geq0$, let
\[
 \Orb_r(x):=
 \bigl\{\alpha\in\Qbar_{>0}:
        \alpha^{x^j}\in\Qbar\text{ for }0\leq j\leq r\bigr\}.
\]
Each $\Orb_r(x)$ is a radical multiplicative subgroup and
\[
 \Orb_0(x)\supseteq\Orb_1(x)\supseteq\Orb_2(x)\supseteq\cdots.
\]
\end{definition}

\begin{theorem}[Orbit filtration]
\label{thm:orbit-filtration}
For a counterexample $x$:
\begin{enumerate}[label=\textup{(\roman*)}]
\item $\Orb_1(x)=\rad_{\Qbar}(\Gammao)$ and has rank $2$;
\item $\Orb_2(x)$ has rank at most $1$;
\item $\Orb_3(x)=\{1\}$, and hence $\Orb_r(x)=\{1\}$ for every $r\geq3$.
\end{enumerate}
Equivalently, no nontrivial positive algebraic number can have four
consecutive algebraic iterates
\[
       \alpha,\quad \alpha^x,\quad
       \alpha^{x^2},\quad \alpha^{x^3}.
\]
\end{theorem}

\begin{proof}
Part (i) is \cref{thm:saturation}.

For (ii), suppose that $\alpha,\beta\in\Orb_2(x)$ are multiplicatively
independent.  Since $x$ is transcendental by \cref{prop:alg-exponent}, the
three numbers $1,x,x^2$ are linearly independent over $\Q$.  The logarithms
$\log\alpha,\log\beta$ are also $\Q$-linearly independent.  All six numbers
\[
 \alpha,\ \beta,\ \alpha^x,\ \beta^x,
 \ \alpha^{x^2},\ \beta^{x^2}
\]
are algebraic, contradicting the Six Exponentials Theorem.  Thus
$\rank\Orb_2(x)\leq1$.

For (iii), assume $\alpha\in\Orb_3(x)$ and $\alpha\neq1$.  Put
\[
 a_0=\alpha,
 \qquad a_1=\alpha^x,
 \qquad a_2=\alpha^{x^2}.
\]
Because $a_j$ and $a_j^x$ are algebraic for $j=0,1,2$, all three elements
$a_0,a_1,a_2$ belong to $\Orb_1(x)$, a group of rank $2$.  They are therefore
multiplicatively dependent: there exist integers $c_0,c_1,c_2$, not all zero,
with
\[
       a_0^{c_0}a_1^{c_1}a_2^{c_2}=1.
\]
Taking real logarithms gives
\[
       \bigl(c_0+c_1x+c_2x^2\bigr)\log\alpha=0.
\]
Since $\alpha\neq1$, this makes $x$ algebraic, contradicting
\cref{prop:alg-exponent}.  Hence $\Orb_3(x)=\{1\}$.
\end{proof}

\begin{corollary}[Forced transcendence at the second or third iterate]
\label{cor:forced-iterate}
For every $q\in\Gammao\setminus\{1\}$, at least one of
\[
             q^{x^2},\qquad q^{x^3}
\]
is transcendental.  In particular, at least one of $2^{x^2}$ and $3^{x^2}$
is transcendental.  If $2^{x^2}$ is algebraic, then $2^{x^3}$ is
transcendental, and similarly with $2$ replaced by $3$.
\end{corollary}

\begin{proof}
Every $q\in\Gammao$ lies in $\Orb_1(x)$.  If both displayed higher iterates
were algebraic, then $q\in\Orb_3(x)$, contradicting
\cref{thm:orbit-filtration}.  If both $2^{x^2}$ and $3^{x^2}$ were algebraic,
then the multiplicatively independent elements $2,3$ would both lie in
$\Orb_2(x)$, contradicting $\rank\Orb_2(x)\leq1$.
\end{proof}

The filtration has an exact interpretation in terms of the multiplicative
rank of the four integers in the problem.  Let
\[
 \rho:=\rank\langle2,3,M,N\rangle.
\]
Since both $\Gammao$ and $\Gammai$ have rank $2$, one has $2\leq\rho\leq4$.

\begin{theorem}[Rank identity and the rank-$3$/rank-$4$ dichotomy]
\label{thm:rank-dichotomy}
For a counterexample,
\[
             \rank\Orb_2(x)=4-\rho.
\]
Consequently $\rho\in\{3,4\}$.

If $\rho=4$, then for every
$\alpha\in\rad_{\Qbar}(\Gammao)\setminus\{1\}$ the number
$\alpha^{x^2}$ is transcendental.

If $\rho=3$, then $\Orb_2(x)$ is a rank-one radical group.  After clearing
radicals there exist integers $a,b,c,d$, with $(a,b)\neq(0,0)$, such that
\[
       (2^a3^b)^x=2^c3^d.
\]
For every nontrivial element $q\in\Orb_2(x)$, the first three iterates
$q,q^x,q^{x^2}$ are algebraic but $q^{x^3}$ is transcendental.
\end{theorem}

\begin{proof}
The power map $\Phi_x$ extends to an isomorphism
\[
  \rad_{\Qbar}(\Gammao)\xrightarrow{\sim}
  \rad_{\Qbar}(\Gammai)
\]
by \cref{thm:saturation}.  Moreover
\[
 \alpha\in\Orb_2(x)
 \quad\Longleftrightarrow\quad
 \alpha\in\rad_{\Qbar}(\Gammao)
 \text{ and }
 \alpha^x\in\rad_{\Qbar}(\Gammao).
\]
Therefore $\Phi_x$ identifies $\Orb_2(x)$ with
\[
 \rad_{\Qbar}(\Gammao)\cap\rad_{\Qbar}(\Gammai).
\]
The rank formula for the product of two finite-rank abelian groups gives
\[
 \rho=2+2-
ank\bigl(
 \rad_{\Qbar}(\Gammao)\cap\rad_{\Qbar}(\Gammai)\bigr),
\]
which is the stated identity.  The bound
$\rank\Orb_2(x)\leq1$ from \cref{thm:orbit-filtration} gives $\rho\geq3$.

When $\rho=4$, the intersection has rank zero; among positive algebraic
numbers its radical torsion is just $1$.  This gives the first case.

When $\rho=3$, the intersection and its preimage have rank one.  Choose a
nontrivial element and raise it to a suitable positive integer power so that
both it and its image lie in the rational lattices $\Gammao$ and $\Gammao$,
respectively.  It then has the displayed form.  The final assertion follows
from $q\in\Orb_2(x)$ and $\Orb_3(x)=\{1\}$.
\end{proof}

Let
\[
        \lambda:=\frac{\log3}{\log2}.
\]
The number $\lambda$ is transcendental: it is irrational, and if it were
algebraic then $2^\lambda=3$ would contradict Gelfond--Schneider.  The return
relation in the rank-$3$ case gives the concrete M\"obius form
\[
       x=\frac{c+d\lambda}{a+b\lambda}.
\]
Thus the rank-$3$ frontier is not an amorphous special case.  It asks whether
a nonconstant rational M\"obius transform of $\lambda$ can simultaneously
send $2$ and $3$ to integers.  The orbit theorem adds a further signature:
the corresponding smooth direction returns once, returns a second time as an
algebraic value, and is forced to become transcendental at the next iterate.

\begin{problem}[Rank-$3$ elimination]
Prove that no nontrivial integers $a,b,c,d$ can satisfy
\[
 (2^a3^b)^x=2^c3^d,
 \qquad 2^x,3^x\in\Z,
 \qquad x\notin\Z.
\]
This would eliminate the entire multiplicative-rank-$3$ branch and leave only
the fully independent rank-$4$ case.
\end{problem}

\section{Rational-span and eigenvalue traps}

The rank dichotomy contains several useful elementary traps.  They are worth
stating separately because they are easy to test against a proposed
counterexample.

\begin{theorem}[Rational-matrix eigenvalue trap]
\label{thm:eigenvalue-trap}
Let
\[
   L=\begin{pmatrix}\log2\\ \log3\end{pmatrix},
   \qquad
   L_1=\begin{pmatrix}\log M\\ \log N\end{pmatrix}=xL.
\]
If either
\[
       L_1=A L
       \quad\text{or}\quad
       L=B L_1
\]
for a $2\times2$ rational matrix $A$ or $B$, then $x\in\Z$.
In particular, if both $M$ and $N$ are $\{2,3\}$-smooth, then $x$ is an
integer.
\end{theorem}

\begin{proof}
In the first case, $AL=xL$, so $x$ is an eigenvalue of a rational matrix and
is therefore algebraic of degree at most two.  In the second case,
$BL_1=L$ gives $BL=x^{-1}L$, so $x^{-1}$ and hence $x$ are algebraic.  Apply
\cref{prop:alg-exponent}.

If $M,N$ are $\{2,3\}$-smooth, their logarithms are integer linear
combinations of $\log2,\log3$, giving the first matrix relation.
\end{proof}

\begin{corollary}[Prime-support and relation-lattice restrictions]
\label{cor:support-rank}
A counterexample has the following properties.
\begin{enumerate}[label=\textup{(\roman*)}]
\item At least one prime $p\geq5$ divides $MN$.
\item The four rational numbers $2,3,M,N$ satisfy at most one independent
multiplicative relation.
\item If there is one relation
\[
       2^a3^bM^cN^d=1,
\]
then necessarily
\[
       x=-\frac{a+b\lambda}{c+d\lambda},
       \qquad \lambda=\frac{\log3}{\log2},
\]
with $c+d\lambda\neq0$.
\end{enumerate}
\end{corollary}

\begin{proof}
If no prime outside $\{2,3\}$ divided $MN$, both outputs would be smooth and
\cref{thm:eigenvalue-trap} would apply.  Two independent relations would make
the multiplicative rank at most two, again impossible by
\cref{thm:rank-dichotomy}.  Taking real logarithms of the single relation and
substituting $\log M=x\log2$, $\log N=x\log3$ gives
\[
      (a+cx)+(b+dx)\lambda=0,
\]
which is the formula.  The denominator cannot vanish because $\lambda$ is
irrational unless $c=d=0$, and then the relation would force $a=b=0$.
\end{proof}

A particularly useful second-iterate diagnostic follows immediately from
saturation:
\[
 \boxed{
  \ 2^{x^2}\in\Qbar
  \iff M\in\rad_{\Qbar}(\Gammao),
  \qquad
  3^{x^2}\in\Qbar
  \iff N\in\rad_{\Qbar}(\Gammao).\ }
\]
Since $M,N$ are rational, ``in the radical'' here means simply
$\{2,3\}$-smooth.  Both conditions cannot hold in a counterexample.

\section{Additive isolation of a hypothetical counterexample}

Multiplication and division stay inside the rank-two domain
$\rad(\Gammao)$, so they cannot produce the third base required by
\cref{cor:one-more-base}.  Addition is the obvious escape route.  The next
result shows exactly how brutally a counterexample must resist it.

\begin{theorem}[Consecutive $\{2,3\}$-smooth integers]
\label{thm:smooth-neighbors}
The positive integers $u$ for which both $u$ and $u+1$ belong to $\smooth$ are
exactly
\[
             u\in\{1,2,3,8\}.
\]
Equivalently, the consecutive smooth pairs are
\[
       (1,2),\quad(2,3),\quad(3,4),\quad(8,9).
\]
\end{theorem}

\begin{proof}
Write
\[
     u=2^a3^b,
     \qquad u+1=2^c3^d,
\]
with nonnegative integer exponents.

If $u$ is odd, then $a=0$.  The case $b=0$ gives $(u,u+1)=(1,2)$.  If
$b\geq1$, then $u+1\equiv1\pmod3$, so $d=0$ and
\[
       2^c-3^b=1.
\]
The elementary edge case gives $2^2-3=1$.  When both exponents exceed one,
Mih\u{a}ilescu's theorem~\cite{Mihailescu} says the only consecutive nontrivial perfect powers
are $8$ and $9$, in the opposite order.  Hence the only odd $u>1$ is $u=3$.

If $u$ is even, then $u+1$ is odd, so $c=0$ and $u+1=3^d$.  If $b\geq1$,
then $u+1\equiv1\pmod3$, impossible for $d\geq1$.  Thus $b=0$ and
\[
       3^d-2^a=1.
\]
The solutions are $3-2=1$ and, by Mih\u{a}ilescu's theorem~\cite{Mihailescu} for the nontrivial
case, $9-8=1$.  Thus $u=2$ or $u=8$.
\end{proof}

\begin{theorem}[Additive isolation]
\label{thm:additive-isolation}
Assume $x$ is a counterexample.  For every $u\in\smooth$,
\[
 (u+1)^x\in\Qbar
 \quad\Longleftrightarrow\quad
 u\in\{1,2,3,8\}.
\]
For all other $u\in\smooth$, the number $(u+1)^x$ is transcendental.
Likewise, for $u>1$ in $\smooth$, $(u-1)^x$ is algebraic exactly for
$u\in\{2,3,4,9\}$.
\end{theorem}

\begin{proof}
By \cref{thm:saturation}, a positive rational $r$ has algebraic $x$-th power
if and only if $r\in\Gammao$.  Apply this criterion to $r=u+1$ and use
\cref{thm:smooth-neighbors}.  The predecessor statement is the same list
shifted by one.
\end{proof}

\begin{corollary}[No long consecutive interpolation block]
The longest block of consecutive positive smooth integers is
\[
              1,2,3,4.
\]
The only other block of length two not contained in it is $8,9$.
Consequently, a finite-difference argument based solely on consecutive
integer values of $t^x$ on $\smooth$ cannot be iterated to arbitrary order.
\end{corollary}

The important point is not merely that $5^x$ is unknown.  Under a
counterexample it is forced to be transcendental.  Thus any prospective
identity that derives algebraicity of $5^x$ (or of the successor of any other
nonexceptional smooth integer) from $M,N\in\Z$ is automatically a complete
proof, not a modest auxiliary step.

\section{A weighted smooth-support series}

The semigroup realization has a natural one-variable encoding:
\begin{equation}
\label{eq:weighted-series}
 H_{M,N}(z):=
 \sum_{a,b\geq0}M^aN^b z^{2^a3^b}\in\Z[[z]].
\end{equation}
Unique factorization guarantees that every exponent occurs at most once.  If
$M=2^x$ and $N=3^x$, the coefficient at a smooth integer $s$ is exactly $s^x$.

Let $T_q$ denote the dilation operator $T_qF(z)=F(z^q)$.

\begin{proposition}[Mixed Mahler identity]
\label{prop:mixed-mahler}
For arbitrary integers $M,N$, not only for a conjectural pair,
\[
       (1-MT_2)(1-NT_3)H_{M,N}(z)=z.
\]
Equivalently,
\[
 H_{M,N}(z)-M H_{M,N}(z^2)-N H_{M,N}(z^3)
       +MN H_{M,N}(z^6)=z.
\]
The two boundary identities are
\[
 H_{M,N}(z)-M H_{M,N}(z^2)=\sum_{b\geq0}N^b z^{3^b},
\]
\[
 H_{M,N}(z)-N H_{M,N}(z^3)=\sum_{a\geq0}M^a z^{2^a}.
\]
\end{proposition}

\begin{proof}
The term $M H(z^2)$ removes precisely the terms with $a\geq1$, while
$N H(z^3)$ removes the terms with $b\geq1$.  Their intersection has been
removed twice and is restored by $MN H(z^6)$.  Only the term $a=b=0$, namely
$z$, remains.  The boundary identities are the corresponding one-direction
versions.
\end{proof}

There is also a Dirichlet-series encoding.  If $c_s$ denotes the coefficient
of $z^s$ in $H_{M,N}$, then
\begin{equation}
\label{eq:dirichlet}
 \sum_{s\geq1}\frac{c_s}{s^w}
  =\sum_{a,b\geq0}\frac{M^aN^b}{(2^a3^b)^w}
  =\frac{1}{(1-M2^{-w})(1-N3^{-w})}.
\end{equation}
The abscissa of absolute convergence is
\[
       \max\{\log_2M,\log_3N\}.
\]
The two quantities are equal exactly when the real dominant singularity is a
double pole.  Under the conjectural relation they both equal $x$, and
\[
 \frac{1}{(1-M2^{-w})(1-N3^{-w})}
  =\frac{1}{(1-2^{-(w-x)})(1-3^{-(w-x)})}.
\]
Thus the original conjecture is equivalently an \emph{integral spectral-shift
problem}: the finite Euler product supported on $\smooth$ is shifted by $x$,
and all shifted Dirichlet coefficients remain integers; prove that the shift
is integral.

\begin{remark}[Why the mixed equation alone cannot prove the conjecture]
The formal identity in \cref{prop:mixed-mahler} holds for every integer pair
$M,N$.  It does not see the archimedean condition
$\log_2M=\log_3N$.  Any argument using only the mixed functional equation and
coefficient integrality is therefore parameter-blind: it would apply equally
to unrelated pairs such as $(M,N)=(5,7)$.  The conjectural information enters
only through coefficient growth, or equivalently through the collision of the
two real pole abscissae in \cref{eq:dirichlet}.
\end{remark}

This is the precise reason that a naive appeal to simultaneous Mahler
rigidity does not work.  We have one equation involving both $T_2$ and $T_3$,
not independent pure-base Mahler equations of the kind to which the standard
Cobham--Mahler rationality principles apply.

\section{Automaticity and quantitative Cartier-state growth}

Let
\[
       F_{2,3}(z):=\sum_{a,b\geq0}z^{2^a3^b}.
\]
Its coefficient sequence is the characteristic sequence of $\smooth$.  We
first give a base-independent obstruction, and then a quantitative binary and
ternary refinement.

\begin{theorem}[The smooth support is automatic in no base]
\label{thm:no-base-automatic}
The set $\smooth=\{2^a3^b:a,b\geq0\}$ is not automatic in any integer base
$k\geq2$.
\end{theorem}

\begin{proof}
Suppose $\smooth$ were $k$-automatic.  Automatic sets are closed under
intersection with ultimately periodic sets.  Intersecting with the integers
not divisible by $3$ gives
\[
       \{2^a:a\geq0\}.
\]
This set is $2$-automatic.  By Cobham's theorem~\cite{Cobham}, if $k$ and $2$ were
multiplicatively independent, it would have to be ultimately periodic, which
it is not.  Therefore $k$ is a power of $2$.

Now intersect $\smooth$ with the odd integers.  This gives
\[
       \{3^b:b\geq0\},
\]
which would be $k$-automatic and is also $3$-automatic.  Since a power of $2$
and $3$ are multiplicatively independent, Cobham's theorem~\cite{Cobham} again forces
ultimate periodicity, a contradiction.
\end{proof}

For a sequence $c=(c_n)_{n\geq0}$ over a finite field, its depth-$r$
$p$-kernel sections are
\[
       c^{(r,s)}=(c_{p^r n+s})_{n\geq0},
       \qquad 0\leq s<p^r.
\]
Equivalently, these are obtained by iterating Cartier operators.  A sequence
is $p$-automatic exactly when its full $p$-kernel is finite.

\begin{theorem}[Exponential Cartier-state growth]
\label{thm:cartier-growth}
Let $c_n$ be the characteristic function of $\smooth$.

For $r\geq3$, the number $K_2(r)$ of distinct binary depth-$r$ kernel sections
satisfies
\[
 K_2(r)\geq
  2^{r-2}-\left\lceil r\log_3 2\right\rceil.
\]
For $r\geq1$, the number $K_3(r)$ of distinct ternary depth-$r$ sections
satisfies
\[
 K_3(r)\geq
  2\cdot3^{r-1}-\left\lceil r\log_2 3\right\rceil.
\]
In particular, $F_{2,3}(z)$ is transcendental over both
$\F_2(z)$ and $\F_3(z)$.
\end{theorem}

\begin{proof}
Fix $r\geq3$ and put $T=\ord_{2^r}(3)=2^{r-2}$.  For
$0\leq j<T$, let $s_j$ be the least nonnegative residue of $3^j$ modulo
$2^r$.  The residue $s_j$ is odd.  Hence a smooth number congruent to $s_j$
modulo $2^r$ must have $2$-adic exponent zero; it is a pure power of $3$.
The relevant exponents are exactly $j+\ell T$, $\ell\geq0$.  Therefore the
first $1$ in the kernel section $c^{(r,s_j)}$ occurs at
\[
       \nu_j=\frac{3^j-s_j}{2^r}
            =\left\lfloor\frac{3^j}{2^r}\right\rfloor.
\]
For
$j\geq J:=\lceil r\log_3 2\rceil$, the values $\nu_j$ are strictly increasing,
because
\[
 \frac{3^{j+1}-3^j}{2^r}
 =\frac{2\cdot3^j}{2^r}\geq2.
\]
Thus the sections indexed by $J\leq j<T$ are pairwise distinct, proving the
binary bound.

For the ternary bound, use
$T=\ord_{3^r}(2)=2\cdot3^{r-1}$ and residues
$s_j\equiv2^j\pmod{3^r}$.  Since these residues are not divisible by $3$, only
pure powers of $2$ occur in the corresponding section.  Its first $1$ is at
\[
       \left\lfloor\frac{2^j}{3^r}\right\rfloor,
\]
and these values are strictly increasing once
$j\geq\lceil r\log_2 3\rceil$.  This proves the ternary bound.

The order identities follow from the standard lifting-the-exponent formulas
\[
 v_2(3^{2^m}-1)=m+2\quad(m\geq1),
 \qquad
 v_3(2^{2\cdot3^m}-1)=m+1\quad(m\geq0).
\]
Finally, Christol's theorem~\cite{ChristolEtAl} identifies algebraicity over $\F_p(z)$ with
$p$-automaticity of the coefficient sequence.  Exponentially many distinct
kernel sections rule out automaticity.
\end{proof}

The next result gives a complete reduction-modulo-prime classification for
the weighted series \cref{eq:weighted-series}.  It is an exact arithmetic
diagnostic, not merely a nonautomaticity statement.

\begin{theorem}[Prime-reduction algebraicity profile]
\label{thm:modp-profile}
Let $\ell$ be a prime, and let $\overline{H}_{M,N,\ell}(z)$ be the reduction of
$H_{M,N}(z)$ modulo $\ell$.  Then
\[
 \overline{H}_{M,N,\ell}(z)
 \text{ is algebraic over }\F_\ell(z)
\]
if and only if one of the following mutually compatible conditions holds:
\begin{enumerate}[label=\textup{(\alph*)}]
\item $\ell\mid M$ and $\ell\mid N$;
\item $\ell=2$ and $2\mid N$;
\item $\ell=3$ and $3\mid M$.
\end{enumerate}
Equivalently,
\[
\begin{array}{c|c}
\ell & \overline{H}_{M,N,\ell}\text{ algebraic}\rule{0pt}{2.4ex}\\
\hline
2 & 2\mid N\\
3 & 3\mid M\\
\ell\geq5 & \ell\mid\gcd(M,N).
\end{array}
\]
\end{theorem}

\begin{proof}
If $\ell\nmid MN$, every coefficient on the smooth support is nonzero.  Were
the coefficient sequence $\ell$-automatic, its nonzero support would also be
$\ell$-automatic, contradicting \cref{thm:no-base-automatic}.

Suppose $\ell\mid M$ but $\ell\nmid N$.  All terms with $a\geq1$ vanish, so
\[
       \overline H(z)=\sum_{b\geq0}\overline N^{\,b}z^{3^b}.
\]
The nonzero weights are periodic in $b$.  Its support is the set of powers of
$3$, which is $3$-automatic and, by Cobham's theorem~\cite{Cobham}, is not
$\ell$-automatic unless $\ell=3$.  When $\ell=3$, the weighted sequence is
indeed $3$-automatic.  Explicitly, if $\eta=\overline N\in\F_3^\times$ and
$K(z)=\sum_{b\geq0}\eta^b z^{3^b}$, then
\[
            K(z)=z+\eta K(z)^3,
\]
so $K$ is algebraic.

The case $\ell\nmid M$ and $\ell\mid N$ is symmetric: only powers of $2$
remain, and algebraicity occurs exactly for $\ell=2$.  In that case
$L(z)=\sum_{a\geq0}z^{2^a}$ satisfies
\[
            L(z)^2+L(z)+z=0
\]
over $\F_2$.

Finally, if $\ell$ divides both $M$ and $N$, every term except $a=b=0$
vanishes and $\overline H(z)=z$.  Combining the cases gives the table.
\end{proof}

\begin{corollary}[What the finite-field profile detects]
The algebraicity profile of the reductions of $H_{M,N}$ recovers:
\begin{itemize}
\item whether $2$ divides the ``wrong'' output $N$;
\item whether $3$ divides the ``wrong'' output $M$;
\item exactly which primes $\ell\geq5$ divide $\gcd(M,N)$.
\end{itemize}
It does not, by itself, detect the equality
$\log_2M=\log_3N$.
\end{corollary}

This last sentence is the key limitation.  The automatic data is highly
nontrivial and exactly computable, but it is local and divisibility-theoretic.
The exponent equation is archimedean.  A proof along this route must couple the
finite-field Cartier profile to the double-pole condition in
\cref{eq:dirichlet}; treating either layer alone leaves arbitrary integer
pairs $M,N$.

\section{Stress tests and discarded shortcuts}

Several attractive-looking shortcuts fail for precise reasons.  Recording
those reasons is useful: it prevents the next round of work from rediscovering
the same dead ends.

\subsection{Six exponentials is saturated, not almost sufficient}

Once $2^x$ and $3^x$ are algebraic, every base in $\Gammao$ has algebraic
$x$-th power.  This produces infinitely many algebraic exponentials, but only
in a two-dimensional multiplicative space.  The Six Exponentials Theorem needs
a third independent logarithmic direction.  More products such as
$6^x=MN$, $12^x=M^2N$, or $(3/2)^x=N/M$ do not increase the rank.

\subsection{The mixed Mahler equation is not two separate Mahler equations}

The identity
\[
 (1-MT_2)(1-NT_3)H=z
\]
does not imply that $H$ satisfies a pure $2$-Mahler equation and a pure
$3$-Mahler equation with rational-function coefficients.  Standard
multiplicatively independent Mahler rigidity cannot therefore be invoked.
The exponential Cartier-state growth in \cref{thm:cartier-growth} is, in fact,
a concrete witness to this obstruction after reduction modulo $2$ or $3$.

\subsection{P\'olya--Carlson behavior is compatible with the series}

The integral power series $H_{M,N}$ is lacunary and normally has the unit
circle as a natural boundary after normalization to radius one.  This is not a
contradiction; it is exactly the alternative allowed by P\'olya--Carlson type
dichotomies.  The mixed functional equation propagates singularities rather
than removing them.

\subsection{Real and $p$-adic exponents cannot be silently identified}

For a prime $p$, one can sometimes solve $2^t=M$ or $3^t=N$ in a $p$-adic
neighborhood after separating torsion and applying the $p$-adic logarithm.
The resulting $p$-adic exponent is not canonically the same object as the real
number $x$.  The real logarithmic determinant gives no reason for the
corresponding $p$-adic determinant to vanish.  Any argument that substitutes a
$p$-adic logarithm for the real one without an explicit global comparison has
inserted the desired conclusion by hand.

\subsection{GCD and LCM do not commute with the power map}

The map $u\mapsto u^x$ preserves products, quotients, order, and divisibility
inside $\Gammao$.  It need not preserve gcd or lcm in the ambient integers,
because $M$ and $N$ may share primes.  Extracting a new base from
$\gcd(M,N)$ would be valuable, but algebraicity of
$\gcd(M,N)^{1/x}$ is not available.

\subsection{Pole alignment is an encoding, not yet a proof}

The double pole of \cref{eq:dirichlet} packages the equality
$\log_2M=\log_3N$ very cleanly.  But a classification theorem saying that an
integral finite Euler product cannot have a nonintegral common pole abscissa
would already be essentially the desired conjecture.  The encoding is useful
only if combined with additional global structure, such as the reduction
profile in \cref{thm:modp-profile}.

\section{Concrete bridge problems}

The new deductions make it possible to formulate sharply targeted next steps.
None of the following is presented as proved.

\begin{problem}[Additive rank raising]
Derive the algebraicity of $(u+1)^x$ for one
$u\in\smooth\setminus\{1,2,3,8\}$ from $M,N\in\Z$.  By
\cref{thm:additive-isolation}, this is immediately contradictory.  The
smallest target is $5^x$, obtained from $u=4$.
\end{problem}

\begin{problem}[Archimedean--Cartier bridge]
Find an invariant of the integral series $H_{M,N}$ that simultaneously sees
\begin{enumerate}[label=\textup{(\roman*)}]
\item the collision of the two real pole abscissae in \cref{eq:dirichlet}, and
\item the prime-reduction algebraicity profile in \cref{thm:modp-profile}.
\end{enumerate}
An invariant depending only on the mixed functional equation cannot work,
because that equation is valid for all $M,N$.
\end{problem}

\begin{problem}[Return-line elimination]
In the rank-$3$ case, eliminate the unique smooth return direction
\[
       (2^a3^b)^x=2^c3^d.
\]
The orbit filtration supplies extra information absent from the original
four-exponential grid: the next iterate is algebraic and the following one is
forced to be transcendental.  A zero estimate or five-of-six exponential
criterion that uses this asymmetric orbit may be enough to rule out the
return line.
\end{problem}

\begin{problem}[Primitive valuation geometry]
Let $u=(v_p(M))_p$ and $v=(v_p(N))_p$ be the two finite-support valuation
vectors, normalized as in \cref{prop:primitive-core}.  Couple their Smith
normal form to the order functional
\[
       (a,b)\longmapsto a\log2+b\log3
\]
and to at least one archimedean small-linear-form estimate.  Pure Smith
normal form is insufficient, but it may provide the missing integral modulus
for an adelic determinant argument.
\end{problem}

Of these, additive rank raising is the cleanest logical target, while the
return-line problem is the narrowest genuinely new transcendence problem.
The automatic/Mahler route remains promising only if it acquires an
archimedean-to-finite-field transfer principle.

\section{Formalization blueprint}

The results split into four nearly independent formalization layers.

\subsection{Elementary arithmetic layer}

This layer requires unique factorization in $\Z$, real logarithms, and basic
finite-rank abelian-group algebra:
\begin{itemize}
\item rational-exponent integrality;
\item multiplicative independence of $M,N$;
\item the ordered-semigroup isomorphism;
\item perfect-power-content descent;
\item the rational-matrix eigenvalue trap;
\item the rank identity
$\rank\Orb_2=4-\rank\langle2,3,M,N\rangle$ once saturation is available.
\end{itemize}

\subsection{Transcendence-theorem interface}

Treat Gelfond--Schneider and Six Exponentials as named imported theorems with
precise hypotheses.  From those interfaces, the proofs of saturation and the
orbit filtration are short.  A Lean-style dependency skeleton is:
\begin{verbatim}
noninteger_solution_transcendental
    -> six_exp_rank_algPowerDomain_le_two
    -> algPowerDomain_eq_radical_two_three
    -> orbit2_rank_le_one
    -> orbit3_eq_one
    -> multiplicative_rank_three_or_four
\end{verbatim}
The delicate formal point is to define multiplicative rank through the
$\Q$-vector space obtained from the torsion-free positive algebraic group.
Restricting to positive reals removes all root-of-unity bookkeeping.

\subsection{Smooth-neighbor layer}

Formalize the parity split in \cref{thm:smooth-neighbors}; the only external
input is Mih\u{a}ilescu's theorem.  The edge cases with exponent $0$ or $1$
should be discharged before invoking the theorem.  Additive isolation then
becomes a one-line corollary of saturation.

\subsection{Automatic-series layer}

The clean sequence of dependencies is:
\begin{verbatim}
Cobham + closure under periodic intersection
    -> smoothSupport_not_automatic_in_any_base
multiplicative_order formulas
    -> quantitative_binary_ternary_kernel_growth
Christol
    -> support_series_transcendental_mod_two_three
case split on prime divisibility of M,N
    -> weighted_reduction_algebraic_iff
\end{verbatim}
The quantitative kernel proof is especially suitable for formalization: each
selected state is distinguished by the position of its first nonzero term.
No minimization over infinite automata is needed.

\section{Conclusions}

The work reported here does not close the Alaoglu--Erd\H{o}s conjecture, but it
substantially sharpens the shape of any counterexample.

A nonintegral solution would have to be transcendental and primitive.  Its
algebraic-power domain would be exactly the rank-two radical generated by
$2,3$; every rational base with any other prime factor would have
transcendental $x$-th power.  The algebraic orbit under repeated exponentiation
would collapse rapidly: the two-step domain has rank at most one, and no
nontrivial algebraic orbit survives through the third iterate.  The four
integers $2,3,M,N$ would have multiplicative rank $3$ or $4$.  Rank $4$ forces
all nontrivial smooth second iterates to be transcendental; rank $3$ permits
one and only one return direction, whose third iterate is then forced to be
transcendental.

On the additive side, the counterexample would be completely isolated.  The
only smooth successors that remain in the algebraic-power domain are the four
classical consecutive smooth pairs ending at $2,3,4,9$.  Any mechanism that
propagates algebraicity as far as $5^x$ proves the conjecture at once.

The weighted support series gives a rigorous computational and formalization
framework.  Its reductions have an exact algebraicity profile and its support
has exponentially many Cartier states in bases $2$ and $3$.  This also shows
why the current Mahler shortcut stops: finite-characteristic automata see
prime divisibility, while the desired conclusion is encoded in a common real
pole abscissa.  The next breakthrough must bridge those two facts or create a
third independent algebraic base.  The rank-$3$ return line and the target
$5^x$ are now the two most concentrated places to attack.

\begin{thebibliography}{99}

\bibitem{AE1944}
L.~Alaoglu and P.~Erd\H{o}s,
\newblock On highly composite and similar numbers,
\newblock \emph{Transactions of the American Mathematical Society}
\textbf{56} (1944), 448--469.

\bibitem{AlloucheShallit}
J.-P.~Allouche and J.~Shallit,
\newblock \emph{Automatic Sequences: Theory, Applications, Generalizations},
\newblock Cambridge University Press, 2003.

\bibitem{ChristolEtAl}
G.~Christol, T.~Kamae, M.~Mend\`es France, and G.~Rauzy,
\newblock Suites alg\'ebriques, automates et substitutions,
\newblock \emph{Bulletin de la Soci\'et\'e Math\'ematique de France}
\textbf{108} (1980), 401--419.

\bibitem{Cobham}
A.~Cobham,
\newblock On the base-dependence of sets of numbers recognizable by finite
automata,
\newblock \emph{Mathematical Systems Theory} \textbf{3} (1969), 186--192.

\bibitem{Gelfond}
A.~O.~Gelfond,
\newblock \emph{Transcendental and Algebraic Numbers},
\newblock Dover, 1960.

\bibitem{Lang}
S.~Lang,
\newblock \emph{Introduction to Transcendental Numbers},
\newblock Addison--Wesley, 1966.

\bibitem{Mihailescu}
P.~Mih\u{a}ilescu,
\newblock Primary cyclotomic units and a proof of Catalan's conjecture,
\newblock \emph{Journal f\"ur die reine und angewandte Mathematik}
\textbf{572} (2004), 167--195.

\bibitem{Schneider}
T.~Schneider,
\newblock \emph{Einf\"uhrung in die transzendenten Zahlen},
\newblock Springer, 1957.

\end{thebibliography}

\end{document}
